\documentclass[../../../../main.tex]{subfiles}
\begin{document}

\begin{flushright}
\footnotesize\textit{【自動文字起こし・要確認】}
\end{flushright}

{\bf [解]}
線分$AB$上の点は
\begin{align*}
\begin{pmatrix}x\\y\\z\end{pmatrix}=s\begin{pmatrix}2\\3\\0\end{pmatrix}+(1-s)\begin{pmatrix}0\\1\\2\end{pmatrix}\quad(0\le s\le1)
\end{align*}
とおける．この時，$k$は$(0\le k\le1)$に対し，
\begin{align*}
s\begin{pmatrix}2\\3\\0\end{pmatrix}+(1-s)\begin{pmatrix}0\\1\\2\end{pmatrix}=k\begin{pmatrix}5+t\\9+2t\\5+3t\end{pmatrix}
\end{align*}
なる$k$があることを示せば良い．
\begin{align*}
2s&=k(5+t)\\
2s+1&=k(9+2t)\\
-2s+2&=k(5+3t)
\end{align*}
\begin{align*}
\therefore\ t=-1，\quad k=\frac13，\quad s=\frac23
\end{align*}

だから，$t=-1$で交点を持ち，その座標は
\begin{align*}
\left(\frac43,\frac73,\frac23\right)．
\end{align*}

\newpage

\end{document}
